% GENERATED FILE. Edit canonical lessons, then run npm run generate:books.
% canonical-source-hash: fnv1a64:d8f41f0ab05d782c
% canonical-lessons: PA-C257-khata, PA-C257-rakam, PA-C257-nakad, PA-C257-kadhvauna, PA-C257-vadhna

\chapter{At the Bank Counter}
\label{ch:pa-a2-at-the-bank-counter}
\hlchaptermodality{257}

\begin{chapteropening}
\textbf{By the end of this chapter:} \emph{I can ask for money from my bank account.}

Complete the last lesson of chapter 257: I can ask for money from my bank account.
\end{chapteropening}

\section[\texorpdfstring{\pa{ਖਾਤਾ} (khātā)}{ਖਾਤਾ (khātā)}]{\pa{ਖਾਤਾ} (khātā) — a bank account; an account book}

\label{lesson:PA-C257-khata}



\begin{quote}
Before the new one: say the Punjabi for to send off, then the Punjabi for to tear.
\end{quote}

\subsection*{You'll want to know: \pa{ਖਾਤਾ}}
\textbf{\pa{ਖਾਤਾ}} — \emph{khātā} — \textquotedblleft{}a bank account; an account book\textquotedblright{}.

\begin{quote}
\pa{ਮੇਰਾ} \pa{ਖਾਤਾ} \pa{ਸ਼ਹਿਰ} \pa{ਦੇ} \pa{ਬੈਂਕ} \pa{ਵਿੱਚ} \pa{ਹੈ।}
\end{quote}

\begin{quote}
\emph{merā khātā shahir de baiṅk vich hai.}
\end{quote}

\begin{quote}
My account is at the bank in town.
\end{quote}

A masculine noun.

\begin{cousinweb}
\textbf{\pa{ਖਾਤਾ}} (\emph{khātā}) is filed by CDIAL under Sanskrit \emph{kṣatra}, \textquotedblleft{}might, rule\textquotedblright{}; CDIAL lists Punjabi \emph{khāttā}, \textquotedblleft{}account book\textquotedblright{}, there, as a word taken from Hindi \emph{khātā} (CDIAL 3648). How the sense moved is not settled.
\end{cousinweb}

\begin{grammarlens}[title={What you've built}]
One more word for reading notices and messages.
\end{grammarlens}

\subsection*{Your turn}
\begin{itemize}
\raggedright
  \item \emph{Say it:} \emph{khātā}
  \item \emph{Say it:} \emph{khātā}, once more
  \item \emph{Say it:} \emph{merā khātā shahir de baiṅk vich hai.} — My account is at the bank in town.
  \item \emph{Recall:} say the Punjabi for to send off, then the Punjabi for to tear, then say \emph{khātā} again
\end{itemize}

\begin{tcolorbox}[breakable,colback=teal!4,colframe=teal!35!black,title={Before you move on}]
What does \pa{ਖਾਤਾ} mean? (\textquotedblleft{}A bank account; an account book\textquotedblright{}.) Say it once more.
\end{tcolorbox}
\section[\texorpdfstring{\pa{ਰਕਮ} (rakam)}{ਰਕਮ (rakam)}]{\pa{ਰਕਮ} (rakam) — a sum of money}

\label{lesson:PA-C257-rakam}



\begin{quote}
Before the new one: say the Punjabi for to tear, then the Punjabi for a bank account.
\end{quote}

\subsection*{You'll want to know: \pa{ਰਕਮ}}
\textbf{\pa{ਰਕਮ}} — \emph{rakam} — \textquotedblleft{}a sum of money\textquotedblright{}.

\begin{quote}
\pa{ਇਹ} \pa{ਰਕਮ} \pa{ਬਹੁਤ} \pa{ਵੱਡੀ} \pa{ਹੈ।}
\end{quote}

\begin{quote}
\emph{eh rakam bahut vaḍḍī hai.}
\end{quote}

\begin{quote}
This sum is very large.
\end{quote}

A feminine noun.

\begin{cousinweb}
\textbf{\pa{ਰਕਮ}} (\emph{rakam}) is a loan, usually traced through Persian to Arabic \emph{raqam}, \textquotedblleft{}a figure, a written sum\textquotedblright{}.
\end{cousinweb}

\begin{grammarlens}[title={What you've built}]
One more word for reading notices and messages.
\end{grammarlens}

\subsection*{Your turn}
\begin{itemize}
\raggedright
  \item \emph{Say it:} \emph{rakam}
  \item \emph{Say it:} \emph{rakam}, once more
  \item \emph{Say it:} \emph{eh rakam bahut vaḍḍī hai.} — This sum is very large.
  \item \emph{Recall:} say the Punjabi for to tear, then the Punjabi for a bank account, then say \emph{rakam} again
  \item \emph{Write it:} copy \textbf{\pa{ਰਕਮ}} once, checking the whole word
\end{itemize}

\begin{tcolorbox}[breakable,colback=teal!4,colframe=teal!35!black,title={Before you move on}]
What does \pa{ਰਕਮ} mean? (\textquotedblleft{}A sum of money\textquotedblright{}.) Say it once more.
\end{tcolorbox}
\section[\texorpdfstring{\pa{ਨਕਦ} (nakad)}{ਨਕਦ (nakad)}]{\pa{ਨਕਦ} (nakad) — in cash}

\label{lesson:PA-C257-nakad}



\begin{quote}
Before the new one: say the Punjabi for a bank account, then the Punjabi for a sum of money.
\end{quote}

\subsection*{You'll want to know: \pa{ਨਕਦ}}
\textbf{\pa{ਨਕਦ}} — \emph{nakad} — \textquotedblleft{}in cash\textquotedblright{}.

\begin{quote}
\pa{ਮੈਂ} \pa{ਦੁਕਾਨ} \pa{ਉੱਤੇ} \pa{ਨਕਦ} \pa{ਪੈਸੇ} \pa{ਦਿੰਦਾ} \pa{ਹਾਂ।}
\end{quote}

\begin{quote}
\emph{maiṁ dukān utte nakad paise dindā hāṁ.}
\end{quote}

\begin{quote}
I pay in cash at the shop.
\end{quote}

\begin{cousinweb}
\textbf{\pa{ਨਕਦ}} (\emph{nakad}) is a loan, usually traced through Persian to Arabic \emph{naqd}, \textquotedblleft{}ready money\textquotedblright{}.
\end{cousinweb}

\begin{grammarlens}[title={What you've built}]
One more word for reading notices and messages.
\end{grammarlens}

\subsection*{Your turn}
\begin{itemize}
\raggedright
  \item \emph{Say it:} \emph{nakad}
  \item \emph{Say it:} \emph{nakad}, once more
  \item \emph{Say it:} \emph{maiṁ dukān utte nakad paise dindā hāṁ.} — I pay in cash at the shop.
  \item \emph{Recall:} say the Punjabi for a bank account, then the Punjabi for a sum of money, then say \emph{nakad} again
  \item \emph{Write it:} copy \textbf{\pa{ਨਕਦ}} once, checking the whole word
\end{itemize}

\begin{tcolorbox}[breakable,colback=teal!4,colframe=teal!35!black,title={Before you move on}]
What does \pa{ਨਕਦ} mean? (\textquotedblleft{}In cash\textquotedblright{}.) Say it once more.
\end{tcolorbox}
\section[\texorpdfstring{\pa{ਕਢਵਾਉਣਾ} (kaḍhvāuṇā)}{ਕਢਵਾਉਣਾ (kaḍhvāuṇā)}]{\pa{ਕਢਵਾਉਣਾ} (kaḍhvāuṇā) — to withdraw; to have something taken out}

\label{lesson:PA-C257-kadhvauna}



\begin{quote}
Before the new one: say the Punjabi for a sum of money, then the Punjabi for in cash.
\end{quote}

\subsection*{You'll want to know: \pa{ਕਢਵਾਉਣਾ}}
\textbf{\pa{ਕਢਵਾਉਣਾ}} — \emph{kaḍhvāuṇā} — \textquotedblleft{}to withdraw; to have something taken out\textquotedblright{}.

\begin{quote}
\pa{ਮੈਂ} \pa{ਹਰ} \pa{ਮਹੀਨੇ} \pa{ਬੈਂਕ} \pa{ਤੋਂ} \pa{ਪੈਸੇ} \pa{ਕਢਵਾਉਂਦਾ} \pa{ਹਾਂ।}
\end{quote}

\begin{quote}
\emph{maiṁ har mahīne baiṅk ton paise kaḍhvāundā hāṁ.}
\end{quote}

\begin{quote}
I withdraw money from the bank every month.
\end{quote}

Past (I …): \pa{ਮੈਂ} \pa{ਬੈਂਕ} \pa{ਤੋਂ} \pa{ਪੈਸੇ} \pa{ਕਢਵਾਏ}, \emph{maiṁ baiṅk ton paise kaḍhvāe}.

\begin{cousinweb}
\textbf{\pa{ਕਢਵਾਉਣਾ}} (\emph{kaḍhvāuṇā}) is the causative of \emph{kaḍḍhṇā}, \textquotedblleft{}to pull out\textquotedblright{}, inherited from a reconstructed \emph{*kaḍḍhati}, \textquotedblleft{}pulls, draws\textquotedblright{}; CDIAL lists Punjabi \emph{kaḍḍhṇā} under it (CDIAL 2660).
\end{cousinweb}

\begin{grammarlens}[title={What you've built}]
One more word for reading notices and messages.
\end{grammarlens}

\subsection*{Your turn}
\begin{itemize}
\raggedright
  \item \emph{Say it:} \emph{kaḍhvāuṇā}
  \item \emph{Say it:} \emph{kaḍhvāuṇā}, once more
  \item \emph{Say it:} \emph{maiṁ har mahīne baiṅk ton paise kaḍhvāundā hāṁ.} — I withdraw money from the bank every month.
  \item \emph{Recall:} say the Punjabi for a sum of money, then the Punjabi for in cash, then say \emph{kaḍhvāuṇā} again
  \item \emph{Write it:} copy \textbf{\pa{ਕਢਵਾਉਣਾ}} once, checking the whole word
\end{itemize}

\begin{tcolorbox}[breakable,colback=teal!4,colframe=teal!35!black,title={Before you move on}]
What does \pa{ਕਢਵਾਉਣਾ} mean? (\textquotedblleft{}To withdraw; to have something taken out\textquotedblright{}.) Say it once more.
\end{tcolorbox}
\section[\texorpdfstring{\pa{ਵਧਣਾ} (vadhṇā)}{ਵਧਣਾ (vadhṇā)}]{\pa{ਵਧਣਾ} (vadhṇā) — to go up; to grow}

\label{lesson:PA-C257-vadhna}



\begin{quote}
Before the new one: say the Punjabi for in cash, then the Punjabi for to withdraw.
\end{quote}

\subsection*{You'll want to know: \pa{ਵਧਣਾ}}
\textbf{\pa{ਵਧਣਾ}} — \emph{vadhṇā} — \textquotedblleft{}to go up; to grow\textquotedblright{}.

\begin{quote}
\pa{ਬੈਂਕ} \pa{ਵਿੱਚ} \pa{ਪੈਸੇ} \pa{ਹੌਲੀ} \pa{ਹੌਲੀ} \pa{ਵਧਦੇ} \pa{ਹਨ।}
\end{quote}

\begin{quote}
\emph{baiṅk vich paise haulī haulī vadhde han.}
\end{quote}

\begin{quote}
Money slowly grows in the bank.
\end{quote}

Past (it …): \pa{ਉਹ} \pa{ਰੁੱਖ} \pa{ਬਹੁਤ} \pa{ਵਧਿਆ}, \emph{oh rukkh bahut vadhiā}.

\begin{cousinweb}
\textbf{\pa{ਵਧਣਾ}} (\emph{vadhṇā}) is inherited from Sanskrit \emph{vardhate}, \textquotedblleft{}grows, increases\textquotedblright{}; CDIAL lists Punjabi \emph{vadhṇā}, \textquotedblleft{}to increase, grow, advance\textquotedblright{}, under it (CDIAL 11376).
\end{cousinweb}

\begin{grammarlens}[title={What you've built}]
One more word for reading notices and messages.
\end{grammarlens}

\subsection*{Your turn}
\begin{itemize}
\raggedright
  \item \emph{Say it:} \emph{vadhṇā}
  \item \emph{Say it:} \emph{vadhṇā}, once more
  \item \emph{Say it:} \emph{baiṅk vich paise haulī haulī vadhde han.} — Money slowly grows in the bank.
  \item \emph{Recall:} say the Punjabi for in cash, then the Punjabi for to withdraw, then say \emph{vadhṇā} again
  \item \emph{Write it:} copy \textbf{\pa{ਵਧਣਾ}} once, checking the whole word
\end{itemize}

\begin{tcolorbox}[breakable,colback=teal!4,colframe=teal!35!black,title={Before you move on}]
What does \pa{ਵਧਣਾ} mean? (\textquotedblleft{}To go up; to grow\textquotedblright{}.) Say it once more.
\end{tcolorbox}
